\documentclass[output=paper,colorlinks=true, citecolor=brown,draftmode]{langscibook}
\ChapterDOI{10.5281/zenodo.18233760}

\author{Nadine Grimm\affiliation{University of Rochester}}
\title{Negation in Gyeli} 
\abstract{The northwestern Bantu language Gyeli exhibits negation features that are generally little discussed explicitly in the Bantu literature regarding the relationship between affirmative constructions and their negative counterparts.
All clausal negation types in Gyeli are asymmetric in \posscitet{Miestamo2005,Miestamo2007} sense. Constructional asymmetries include tone changes and changes in finiteness, while paradigmatic asymmetries mostly concern TAM neutralization; aspect marking is generally lost under negation (unless it is specifically marked in an embedded clause). The negation pattern in the present tense is distinct from other tense-mood categories in its morphological negation-marking strategy and a higher degree of asymmetries. The present tense-mood category also allows a choice between different negation markers that are linked to information structure and focus marking.

\keywords{TAM systems, tonal morphology, Bantu, negation, information structure}}

\IfFileExists{../localcommands.tex}{
   \addbibresource{../localbibliography.bib}
   % add all extra packages you need to load to this file

\usepackage{tabularx,multicol}
\usepackage{url}
\urlstyle{same}

\usepackage{listings}
\lstset{basicstyle=\ttfamily,tabsize=2,breaklines=true}

% \usepackage{langsci-basic}
\usepackage{langsci-optional}
\usepackage{langsci-lgr}
\usepackage{langsci-gb4e}

\usepackage[linguistics,edges]{forest}
\usepackage{subfigure}

\usepackage{arydshln}

\usepackage{pifont}
\newcommand{\cmark}{\ding{51}}

\usepackage{tabto}


\usepackage{multirow} %for Tsova-Tush, NGT, panara
\usepackage{booktabs} %for Tsova-Tush, panara
\usepackage{graphicx} %for Tsova-Tush
\usepackage{longtable} %for Tsova-Tush, NGT
\usepackage{vcell} %for Tsova-Tush

% \usepackage{pdflscape} %kambaata -- has been added to enable landscape formatting of table 3
\usepackage{xcolor} %kambaata,ngarinjin,nivkh -- has been added to enable gray shading in tables
\definecolor{light-gray}{gray}{0.8}%ngarinjin
\definecolor{lightest-gray}{gray}{0.95}%ngarinjin

%\usepackage{subcaption} %for NGT, but does not work with subfigure package from above
\usepackage{slgloss} %for NGT
% \newfontfamily{\handshape}{handshape.TTF} %for NGT %HC2022-11-04: it does not work and creates an error "Package fontspec Error: The font "handshape" cannot be found", eventhough I have uploaded the file handshape.TTF, so I outcommented it for the time being. Could you solve this Sebastian? thanks :)
%\newcommand\customfont[1]{{\usefont{T1}{handshape}{m}{n} #1 }} %for NGT
\newfontfamily\ToneFont{CharisSIL-Regular.ttf}
\AdditionalFontImprint{Charis SIL}
\usepackage[shortlabels]{enumitem} %NGT 


\usepackage{leipzig} %for the glosses in panara
\usepackage{glossaries} %panara
% \usepackage{lscape} %panara
% \input{libertine-preamble} %panara
% \setmathfont{Libertinus Math} %panara
%   \usepackage{enumerate} %panara - it clashes with enumitem which is really needed for NGT, and it works in panara also with enumitem, i.e. we don't need enumerate.

% \setmainfont[Ligatures=TeX,Numbers=OldStyle]{Linux Libertine O} %panara
% \setsansfont[Ligatures=TeX,Numbers=OldStyle]{Linux Biolinum O} %panara

\usepackage{hyperref}%kalamang,yukana
% \usepackage{tipx}%kalamang,ungarinjin, yukana

% \usepackage{wrapfig}%ungarinjin
% \usepackage{tabulary}%ungarinjin
%\usepackage{subfig}%ungarinjin %deleted bc subfigure already above
\PassOptionsToPackage{colorinlistoftodos,disable}{todonotes}
\usepackage{todonotes}
\usepackage{hyphenat}
\usepackage{langsci-branding}


   %\newcommand*{\orcid}{}

\makeatletter
\let\thetitle\@title
\let\theauthor\@author
\makeatother

\newcommand{\togglepaper}[1][0]{
   \bibliography{../localbibliography}
   \papernote{\scriptsize\normalfont
     \theauthor.
     \titleTemp.
     To appear in:
     Matti Miestamo \& Ljuba Veselinova (ed.).
     Title of book.
     Berlin: Language Science Press. [preliminary page numbering]
   }
   \pagenumbering{roman}
   \setcounter{chapter}{#1}
   \addtocounter{chapter}{-1}
}


\renewbibmacro*{index:name}[5]{% %Tsova-Tush
  \usebibmacro{index:entry}{#1} %Tsova-Tush
    {\iffieldundef{usera}{}{\thefield{usera}\actualoperator}\mkbibindexname{#2}{#3}{#4}{#5}}} %for Tsova-Tush

    
\newcommand{\gl}[1]{\textsc{#1}}

\newcommand{\ipa}[1]{{\textit{#1}}} %geshiza 
\newcommand{\ipaEX}[1]{{\textit{#1}}} %geshiza


%\newcommand{\hspaceThis}[1]{\hphantom{#1}} %this was in localcommands in the kambaata chapter

\newcommand{\ANA}{\textsc{ana}} %gyeli
\newcommand{\ATTR}{\textsc{attr}} %gyeli
\newcommand{\CONJ}{\textsc{conj}} %gyeli
\newcommand{\ID}{\textsc{id}} %gyeli
\newcommand{\IDEO}{\textsc{ideo}} %gyeli
\newcommand{\INCH}{\textsc{inch}} %gyeli
\newcommand{\LINK}{\textsc{link}} %gyeli
\newcommand{\NP}{NP} %gyeli
\newcommand{\NCA}{\textsc{nca}} %gyeli
\newcommand{\NPI}{\textsc{npi}} %gyeli
\newcommand{\PCF}{\textsc{pcf}} %gyeli
\newcommand{\PN}{\textsc{pn}} %gyeli
\newcommand{\PREP}{\textsc{prep}} %gyeli
\newcommand{\PRIOR}{\textsc{prior}} %gyeli
\newcommand{\PRO}{\textsc{pro}} %gyeli
\newcommand{\PROSP}{\textsc{prosp}} %gyeli
\newcommand{\R}{\textsc{r}} %gyeli
\newcommand{\RETRO}{\textsc{retro}} %gyeli
\newcommand{\SEQU}{\textsc{sequ}} %gyeli
\newcommand{\STAMP}{\textsc{stamp}} %gyeli
\newcommand{\SUB}{\textsc{sub}} %gyeli
\newcommand{\TM}{TM} %gyeli
\newcommand{\V}{\textsc{v}} %gyeli


\newcommand{\hspaceThis}[1]{\hphantom{#1}} %for Kambaata chapter, helps for \db (which allows alignment of word and gloss and not without taking into account the "s)


%for Panara:
\renewbibmacro*{index:name}[5]{%
  \usebibmacro{index:entry}{#1}
    {\iffieldundef{usera}{}{\thefield{usera}\actualoperator}\mkbibindexname{#2}{#3}{#4}{#5}}}

 %panara
\newleipzig{ade}{ade}{adessive} %panara
\newleipzig{dir}{dir}{directional} %panara
\newleipzig{emph}{emph}{emphasis} %panara
\newleipzig{fnl}{fnl}{final} %panara
\newleipzig{nml}{nml}{nominal form}%panara
\newleipzig{ine}{ine}{inessive}  %panara 
\newleipzig{ints}{ints}{intensifier} %panara
\newleipzig{iter}{iter}{iterative} %panara
\newleipzig{plac}{plac}{pluractional} %panara


% \makeglossaries %panara %HC not sure what to do with these following lines, but it does not work yet. I have added the abbreviations in the same table is in other chapters, so this glossary function is not needed anymore.
% \newglossarystyle{myglosses}{%
%   \renewenvironment{theglossary}%
% {\begin{multicols}{2}\raggedright}
% {\end{multicols}}
%     \renewcommand*{\glossaryheader}{}
%     \renewcommand*{\glsgroupheading}[1]{}
%     \renewcommand*{\glsgroupskip}{}
%     \renewcommand*{\glsclearpage}{} 
% % set how each entry should appear:
%     \renewcommand*{\glossentry}[2]{
%     \noindent\makebox[7em][l]{\glstarget{##1}{\textsc{\glossentryname{##1}}}}
%         \glossentrydesc{##1}
%     }
%     \renewcommand*{\subglossentry}[3]{%
%         \glossentry{##2}{##3}
%     }
% }

\newcommand{\glopa}{\textsc{dem}} %kalamang
\newcommand{\glme}{\textsc{top}} %kalamang
\newcommand{\glte}{\textsc{nfin}} %kalamang
\newcommand{\glse}{\textsc{iam}} %kalamang
\newcommand{\glkin}{\textsc{vol}} %kalamang
\newcommand{\glteba}{\textsc{prog}} %kalamang
\newcommand{\gli}{\textsc{plnk}} %kalamang
\newcommand{\glet}{\textsc{irr}} %kalamang
\newcommand{\glta}{\textit{ta}} %kalamang
\newcommand{\glnn}{\textit{n}} %kalamang

\newcommand{\tabitem}{~~\llap{\textbullet}~~}%kogi

%For Matti and Ljuba's comments on the last version of the chapters:

\newcommand{\matti}[2][color=green!40,size=\small]{\todo[#1]{MM: #2}}
\newcommand{\ljuba}[2][color=blue!40,size=\small]{\todo[#1]{LV: #2}}
\newcommand{\mm}[2][color=green!40,size=\small]{\todo[#1]{MM: #2}}
\renewcommand{\lv}[2][color=blue!40,size=\small]{\todo[#1]{LV: #2}}
\newcommand{\hc}[2][size=\small]{\todo[#1]{HC: #2}}
\newcommand{\gloss}[2][color=red!30,size=\tiny]{\todo[#1]{Glossing: #2}}
\newcommand{\glosst}[2][color=red!30,size=\tiny,inline]{\todo[#1]{Glossing: #2}} %t for (in the) text / inline
\newcommand{\aut}[2][color=yellow!50,size=\tiny]{\todo[#1]{Author: #2}}%for author

\renewcommand{\keywords}[1]{%
  \par\noindent\textbf{Keywords: #1.}%
}


\renewcommand{\REF}[2][]{(\ref{#2}#1)}
\renewcommand{\sectref}[1]{§\ref{#1}}
\newcommand{\Abf}{{Ā}\kern-2pt\raisebox{3.1pt}{́}\kern2pt}
\newcommand{\abf}{{ā}\kern-1pt\raisebox{1.1pt}{́}\kern1pt}
\newcommand{\A}{{Ā}\kern-1.5pt\raisebox{3.1pt}{́}\kern1.5pt}
\newcommand{\U}{{ʊ}\kern-1.5pt\raisebox{-.4pt}{̀}\kern1.5pt}

\newcommand{\fdot}{{f\kern-.75pt\raisebox{-2.5pt}{̣}\kern.75pt}}

\newcommand{\handshape}[2][1.5]{\includegraphics[height=#1ex]{fonts/#2.png}}

    \togglepaper[23]
}{}

\begin{document}
\maketitle

\section{The language}

Gyeli (ISO: \href{https://iso639-3.sil.org/code/gyi}{gyi}, Glottocode: \href{https://glottolog.org/resource/languoid/id/gyel1242}{gyel1242}) is an endangered Bantu A80 language spoken by so-called ``Pygmy'' hunter-gatherers in the rainforest of southern Cameroon. The language is also known under the names Gyele, Kola, Bagielli, or Bakola, while the speakers call themselves Bagyeli or Bakola depending on the region where they live: in the northern part, the term Bakola is used, in the central and southern part Bagyeli. Speaker numbers are estimated to range between 4,000 and 5,000 ({\itshape Ethnologue}, \citealt{Ngima2001}). The speech community is dispersed over a vast area of about 12,500 km\textsuperscript{2} which they share with eight other ethnic groups: Batanga, Bakoko, Basaa, Ewondo, Bulu, Fang, Yasa, and Kwasio. These groups all speak other Bantu A languages and are all agriculturalists and/or fishermen (rather than forest foragers).

Speakers are currently changing their traditional way of life due to massive changes in the environment. In 2015, the biggest deep sea port for central Africa was inaugurated just off the shore where the Bagyeli live, resulting in a significant increase in infrastructure development which leads to a loss of rainforest and as such loss of the animals that the Bagyeli as hunter-gatherers depend on. As a consequence, the Bagyeli are being forced to give up hunting and rely more on farming just as their neighboring communities do. In the course of this subsistence change, the Bagyeli are also shifting to the languages of other Bantu groups. This shift is reinforced by the low prestige of the language and the marginalization of the Bagyeli within Cameroonian society at large. Different dialects of Gyeli are currently emerging, depending on which language a particular Gyeli group is mostly in contact with. There is no standard variety of Gyeli, which correlates with the lack of an official written language. Instead, Gyeli is currently used orally only. The conventions in \posscitet{Grimm2021} description, however, have been informed by speakers' preferences and are potentially a first step towards a written form of Gyeli.  In general, everybody in this area, both Bagyeli and other Bantu groups, is multilingual and speaks easily up to five languages. In contrast to the farming Bantu groups of Cameroon, the Bagyeli generally do not speak the colonial language French or speak it less fluently, since Bagyeli children receive less or no schooling.  

The main work that exists on the Gyeli language consists of a description of the phonology and nominal morphology in the Kwasio and Basaa contact region by \citet{Renaud76} and \posscitet{Grimm2021} reference grammar of the Gyeli variety spoken in the village Ngolo in the Bulu contact region. A documentary corpus of Gyeli is available in {\itshape The Language Archive}, including natural texts and elicitations \citep{Grimm2020}.  Data for this paper were obtained during fieldwork in Cameroon between 2010 and 2017 as part of the language documentation and grammar-writing project. A detailed description of consulted speakers and the various methodologies used is provided in \citet[28--29]{Grimm2021}.

Linguistically, Gyeli is an S V O (X), head-initial language. It features a gender system with nine agreement classes and six genders.\footnote{In the examples in this chapter, only agreement classes are indicated in glossing nominal prefixes. In \citet{Grimm2021}, I also indicate noun prefix classes in my glosses.}
In comparison to eastern and southern Bantu languages, Gyeli is more analytic and has, for instance, a limit of three syllables in verb stems. Also, Gyeli is a tone language, phonologically distinguishing H, L, HL, LH, and toneless tone-bearing units, which encode both lexical and grammatical distinctions. 

In this paper, following the questionnaire provided by the editors (\citetv{chapters/questionnaire}),
% \mm{@LSP: cross-ref to the questionnaire}
I first discuss clausal negation in \sectref{sec:gye:ClausNEG}, distinguishing standard negation from non-standard negation marking. In \sectref{sec:gye:3}, I describe aspects of non-clausal negation, including negative replies and the absence of negative indefinites, quantifiers, and derivation. In \sectref{sec:gye:4}, I address other aspects of negation, such as the scope of negation, negative polarity, and negation in complex clauses. Finally, \sectref{sec:gye:5} concludes the chapter.

\section{Clausal negation}
\label{sec:gye:ClausNEG}

Following \citet[42]{Miestamo2005} and later accounts of negation typology \citep{AuweraKrasnoukhova2020}, I distinguish standard from non-standard negation strategies.  \citet[42]{Miestamo2005} defines standard  negation  as ``the productive and general means of negating declarative verbal main clauses''. 

The remainder of this section is structured as follows: in \sectref{sec:gye:SN}, I describe standard negation with subsections on constructional and paradigmatic asymmetries between positive verbal clauses and their negative counterparts. Negation in non-declaratives, including questions, imperatives, and subjunctives, is described in \sectref{sec:gye:SNnondecl}. The negation of stative predications is discussed in \sectref{sec:gye:Stative}, while non-main clauses and their negation strategies are explored in \sectref{sec:gye:non-main}.  

\subsection{Standard negation}
\label{sec:gye:SN}

In this section, I first introduce the negative markers (\sectref{sec:gye:NegMarkers}). In \sectref{sec:gye:AIntro}, I introduce negation asymmetries, presenting constructional asymmetries in detail in \sectref{sec:gye:CA} and paradigmatic asymmetries in \sectref{sec:gye:PA}. 

\subsubsection{Negative markers}
\label{sec:gye:NegMarkers}

Gyeli has three standard negation markers, each for different tense-mood categories, as shown in \tabref{tab:2:negmarkers1}. Regarding the typology of morphological versus syntactic negation marking \citep{Dahl1979,Payne1985}, most negation markers in Gyeli are auxiliaries and thus of the syntactic type (see \tabref{tab:1:negmarkers}). The only morphological negation marker is found in present tense negation. This marker is a circumfix, consisting of a segmental suffix -{\itshape lɛ} and a H tone co-exponent that attaches to the left of the verb, spreading rightwards. The negation auxiliaries are morphologically complex with a final syllable -{\itshape lɛ} that seems related to the present negation. Synchronically, however, the auxiliaries cannot be parsed into lexical and negation morphemes. I therefore consider them monomorphemic.

\begin{table}
    \caption{Markers of standard clausal negation in Gyeli}
    \label{tab:2:negmarkers1}
    \begin{tabular}{lll}
      \lsptoprule
        Negation marker & Status & Distribution \\
        \midrule
        \multicolumn{3}{l}{Standard negation} \\
        {\bfseries H...-lɛ} &  negation circumfix & present \\
        {\bfseries sàlɛ́/pálɛ́} &  negation auxiliary & past tenses \\
        {\bfseries kálɛ̀} &  negation auxiliary & future \\
      \lspbottomrule
    \end{tabular}
\end{table}

Examples for each negation marker and its affirmative counterpart are given in (\ref{SN1}--\ref{SN3}).\footnote{Further example sentences in \sectref{sec:gye:CA} and \sectref{sec:gye:PA} illustrate the different types of asymmetries.} The tone rules are explained in \sectref{sec:gye:AIntro}. In a nutshell, tone patterns on the preverbal clitic and the verb differ depending on the tense-mood category they encode as well as on a phrase-medial versus phrase-final position of the verb and polarity. They are not conditioned by the lexical tone of the verb. 

\ea\label{SN1}
    \ea\label{SN1x}
        \glll bá dé bédéwɔ̀ \\ 
              ba-H dè-H {H\backslash{}be-déwɔ̀} \\
              2-{\PRS} eat-{\R} {{\OBJ}.{\LINK}\backslash{}8-food} \\
        \trans `They eat food.'
    \ex\label{SN2x}
    \glll    bá dé{\bfseries lɛ́} bédéwɔ̀ \\
              ba-H dè-lɛ {H\backslash{}be-déwɔ̀}\\
                 2-{\PRS} eat-{\NEG} {{\OBJ}.{\LINK}\backslash{}8-food} \\
        \trans `They don't eat food.'
    \z
\z

\ea\label{SN2}
    \ea\label{SN2xa}
    \glll    bà dé bédéwɔ̀ \\
              ba dè-H {H\backslash{}be-déwɔ̀}\\
        2.{\textsc{rec.}\PST} eat-{\R} {{\OBJ}.{\LINK}\backslash{}8-food} \\
        \trans `They ate food.'
     \ex\label{SN2xb}
    \glll    bà {\bfseries sàlɛ́/pálɛ́} dè bédéwɔ̀ \\
              ba sàlɛ́/pálɛ́ dè {H\backslash{}be-déwɔ̀}\\
                 2.{\textsc{rec.}\PST} {\NEG}.{\PST} eat {{\OBJ}.{\LINK}\backslash{}8-food}\\
        \trans `They didn't eat food.'
    \z    
\z

\ea\label{SN3}    
    \ea\label{SN3xa}
    \glll    báà dè bédéwɔ̀ \\
              báà dè {H\backslash{}be-déwɔ̀}\\
                 2.{\FUT} eat {{\OBJ}.{\LINK}\backslash{}8-food}\\
        \trans `They will eat food.' 
    \ex\label{SN3xb}
    \glll    báà {\bfseries kálɛ̀} dè bédéwɔ̀ \\
              báà kálɛ̀ dè {H\backslash{}be-déwɔ̀}\\
                 2.{\FUT} {\NEG}.{\FUT} eat {{\OBJ}.{\LINK}\backslash{}8-food}\\
        \trans `They won't eat food.'    
    \z    
\z

\subsubsection{Introducing negation asymmetries}
\label{sec:gye:AIntro}

Gyeli negation is asymmetric in \posscitet{Miestamo2007} sense with structural differences between affirmatives and their negative counterparts, which mostly play out in the tonal patterns of the verb stem and the preverbal portmanteau morpheme ``{\STAMP}'', which encodes subject agreement, tense, aspect, mood, and polarity. These differences are both constructional and paradigmatic, as discussed in detail below. 

In order to follow my argumentation on negation asymmetries, it is important to understand how Gyeli encodes tense categories in affirmative basic verbal predicates. The primary marking of tense in affirmative clauses is tonal, unlike negation marking, which is achieved primarily through segmental material (morphological and syntactic markers). The combinations of tonal patterns on the {\STAMP} morpheme, a preverbal clitic, and the verb stem instantiate seven distinct categories, as shown in \tabref{Tab:TM-Tone}.

\begin{table}
    \caption{Tonal patterns of affirmative tense-mood categories}
    \label{Tab:TM-Tone}
    \small 
    \begin{tabularx}{\textwidth}{lllllX}
        \lsptoprule
        Basic  & {\TM}  & \textsc{stamp} & Verb& Verb & Gloss\\
        distinction &  category &  & Stem & Tone &   \\
        \midrule
        & \textsc{prs} & {\itshape yá} & {\itshape dè} & \multirow{3}{*}{L} & `We eat.' \\
        non-past & \textsc{inch} & {\itshape yàá} & {\itshape dè} & & `We are at the beginning of eating.' \\
        & \textsc{fut} & {\itshape yáà/mɛ̀ɛ̀} &  {\itshape dè}  & &  `We/I will eat.' \\
        \midrule
        \multirow{2}{*}{past} &  \textsc{rec.pst} & {\itshape yà} &  {\itshape dé} & \multirow{2}{*}{H} &  `We ate (recently).' \\
        & \textsc{rem.pst} & {\itshape yáà} & {\itshape dé} & &  `We ate (a long time ago).' \\
        \midrule
        \multirow{2}{*}{tenseless} & \textsc{imp} &  ({\itshape yá})  & {\itshape dê}    & \multirow{2}{*}{HL} &  `Let's eat!' \\
        & \textsc{sbjv} & {\itshape yá}  & {\itshape déè} & &  `May we eat.' \\
        \lspbottomrule
    \end{tabularx}
\end{table}

Tense is intrinsically linked to mood marking, distinguishing realis and irrealis. Each tense inherently belongs to either the realis or the irrealis mood.  I refer to these categories as ``tense-mood'' ({\TM}) categories. Mood is marked overtly only for the realis category by an H tone that attaches to the right of the verb stem -- if the verb is not phrase-final. If the verb is phrase-final, mood marking tones do not surface, as I show below.  In fact, the only inflection marking on finite verbs is mood marking opposing realis (with H) and irrealis (no H) categories. 

The presence of the realis marking H tone is conditioned by the specific {\TM} category and not morphosyntactic material following the verb since all parts of speech following the verb trigger the syntactic tone to surface.
This is the case, for instance, in present affirmative clauses, as in \REF{MetaPOS}, where a range of different word classes that follow the finite verb trigger a final H tone on the verb. I gloss this H tone as {\R} for realis. 

\eabox{\label{MetaPOS}
\begin{tabular}{llll}
a. & mɛ́ gyámbɔ̀                   &  `I cook.'                 & \\
b. & mɛ́ gyámbɔ́ bé-kwàndɔ̀         &  `I cook plantains.'       &  $\underline{\quad}${\N} \\
c. & mɛ́ gyámbɔ́ byɔ̂               &  `I cook it.'              & $\underline{\quad}${\PRO} \\
d. & mɛ́ gyámbɔ́ ndáà              &   `I cook today.'          &  $\underline{\quad}${\ADV} \\
e. & mɛ́ gyámbɔ́ ɛ́ kìsíní dé tù    &   `I cook in the kitchen.' &  $\underline{\quad}${\PREP} \\
f. & mɛ́ gyámbɔ́ nà wɔ́mbɛ̀lɛ̀        &  `I cook and sweep.'       &  $\underline{\quad}${\CONJ} \\
g. & mɛ́ wúmbɛ́ gyámbɔ̀             &  `They want to cook.'      &  $\underline{\quad}${\V} \\
\end{tabular}
}
\tabref{tab:3:moods} shows which affirmative and negative tense categories are  linked to which mood category. Thus, basic verbal clauses in the present, inchoative, recent, and remote past as realis categories always surface with a final H on the finite verb, while finite verbs in the irrealis categories of future, imperative, and subjunctive end in a non-H tone. The surface tones of H versus non-high clearly partition the tense-mood categories into realis/irrealis, constituting a semantic pattern that should not be easily dismissed as an arbitrary phonological pattern. Past negative auxiliaries {\itshape sàlɛ́} and {\itshape pálɛ́} are in the realis category, just like their affirmative counterparts. Mood category alignment is also found with the irrealis future in both affirmative and negative. Only the present switches from a realis category in the affirmative to irrealis under negation (see \sectref{sec:gye:PA} on paradigmatic asymmetries). 

\begin{table}
    \caption{Distribution of realis and irrealis categories}
    \label{tab:3:moods}
    \begin{tabular}{lll}
        \lsptoprule
         & H tone presence & H tone absence\\
         & $\rightarrow$ Realis & $\rightarrow$ Irrealis\\
         \midrule
        \multirow{4}{*}{Affirmative} & present & future \\
        & inchoative & imperative  \\
        & recent past  & subjunctive \\
        & remote past &  \\
        \midrule 
        \multirow{2}{*}{Negative} & past negatives & present\\
         &  & future negative \\
         \lspbottomrule
    \end{tabular}
\end{table}

Having outlined the tense-mood marking system of basic verbal predicates in affirmatives, I now turn to describing the asymmetries that occur under negation. 

\subsubsection{Constructional asymmetries}
\label{sec:gye:CA}

Standard clausal negation is constructionally asymmetric with both the negation circumfix H-...-{\itshape lɛ} and the negation auxiliaries {\itshape sàlɛ/pálɛ́} and {\itshape kálɛ̀}, although the nature of the asymmetry differs between the morphological and the syntactic negation strategies. 

With the negation circumfix H...-{\itshape lɛ} in the present tense-mood category, the {\STAMP} changes its tonal pattern (and vowel length) under negation, but only for the first and second person singular and agreement class 1. These person markers receive a long vowel with a rising tone, as shown for the first person singular in \REF{CA}. 

\ea\label{CA}
\ea\label{CA3}
\glll    {\bfseries mɛ́} dé bédéwɔ̀ \\ 
          mɛ-H dè-H {H\backslash{}be-déwɔ̀}\\
             1{\SG}-{\PRS} eat-{\R} {{\OBJ}.{\LINK}\backslash{}8-food}\\
    \trans `I eat food.'
     \ex\label{CA4}
\glll    {\bfseries mɛ̀ɛ́} délɛ́ bédéwɔ̀ \\
          mɛ̀ɛ́ dè-lɛ {H\backslash{}be-déwɔ̀}\\
             1{\SG}.{\PRS}.{\NEG} eat-{\NEG} {{\OBJ}.{\LINK}\backslash{}8-food}\\
    \trans `I don't eat food.'
\z
\z

All other {\STAMP} morphemes (first and second person plural and agreement classes 2 through 9) are not subject to this asymmetry and retain their H tone pattern on the {\STAMP} under negation, as shown with agreement class 2 in \REF{SN1}. In these cases, negation is symmetric. 

The asymmetry of negation auxiliaries {\itshape sàlɛ́/pálɛ́} in past tenses and {\itshape kálɛ̀} in the future concerns a change in finiteness. This type of asymmetry is what \citet{Miestamo2007} calls ``A/Fin/NegVerb''. In affirmative simple verbal predicates,  the lexical verb is also the finite verb, as in \REF{meta2a}. In contrast, under negation, the negation auxiliary becomes the finite element (as seen also above in \ref{SN2} and \ref{SN3}), while the lexical verb is not inflected for realis marking, as shown in \REF{meta2b} for the recent past. The remote past and the future are structurally identical to \REF{meta2b}.

\ea\label{meta2}
\ea\label{meta2a} 
  \glll     mɛ̀ gyám{\bfseries bɔ́} bélɔ̀lɔ̀ \\
          mɛ gyámbɔ-H {H\backslash{}be-lɔ̀lɔ }\\
              1{\SG}.{\textsc{rec.}\PST} cook-{\R} {{\OBJ}.{\LINK}\backslash{}8-duck}\\
    \trans `I cooked ducks.'
\ex\label{meta2b}
  \glll     mɛ̀ sàlɛ́  gyám{\bfseries bɔ̀} bélɔ̀lɔ̀\\
            mɛ sàlɛ́  gyámbɔ {H\backslash{}be-lɔ̀lɔ }\\
              1{\SG}.{\PST} {\NEG}.{\PST}.{\R}  cook {{\OBJ}.{\LINK}\backslash{}8-duck}\\
    \trans `I did not cook ducks.'
\ex\label{meta2c}
  \glll     mɛ́ wúm{\bfseries bɛ́}  gyám{\bfseries bɔ̀} bélɔ̀lɔ̀ \\
            mɛ-H wúmbɛ-H  gyámbɔ {H\backslash{}be-lɔ̀lɔ }\\
              1{\SG}-{\PRS} want-{\R} cook {{\OBJ}.{\LINK}\backslash{}8-duck}\\
    \trans `I want to cook ducks.'
\z 
\z

As negation auxiliaries never occur in isolation, i.e. outside of a complex predicate, their underlying tonal structure cannot be proven. I still argue that negation auxiliaries become the finite verbal element in complex predicates based on their parallel structure to modal auxiliaries such as {\itshape wúmbɛ} `want' in \REF{meta2c}. Modal auxiliaries serve as the finite element in complex predicates, while the realis marking H tone is consistently absent on the lexical verb {\itshape gyámbɔ̀} `cook'.

\subsubsection{Paradigmatic asymmetries}
\label{sec:gye:PA}

All standard negation strategies in verbal clauses are also subject to paradigmatic asymmetries. First, the present tense switches from a realis mood category in affirmative clauses to an irrealis category under negation with the circumfix H...\nobreakdash-{\itshape lɛ}. Example \REF{leTone} shows that the realis marking H tone is present in the affirmative clause, but absent in its negative counterpart. 

\ea\label{leTone}
\ea\label{leTonea} 
  \glll     Àdà á gyà{\bfseries gá} mábèlé \\
          Àdà a-H gyàga-H {H\backslash{}ma-bèlé}\\
          $\emptyset$.1.PN 1-{\PRS} buy-{\R} {{\OBJ}.{\LINK}\backslash{}6-kola.nut}\\
    \trans `Ada buys kola nuts.'
\ex\label{leToneb}
  \glll    Àdà  àá {\bfseries gyágàlɛ̀} mábèlé\\
          Àdà àá H-gyàga-lɛ {H\backslash{}ma-bèlé}\\
      $\emptyset$.1.PN 1.{\PRS}.{\NEG} {\NEG}-buy-{\NEG} {{\OBJ}.{\LINK}\backslash{}6-kola.nut}\\
    \trans `Ada does not buy kola nuts.'
\z 
\z

Instead of simply lacking the realis marking H tone, however, the tone assignment rules for the present negation are different than for the affirmative realis/irrealis distinction. First, the \textsc{stamp} marker takes a different form in the first and second person singular and in agreement class 1 with a long LH vowel instead of a plain H morpheme, as shown in \sectref{sec:gye:CA} for a constructional asymmetry. Second, the surface tone of the verb is determined by the tonal co-exponent of the present negation circumfix /H-...-lɛ/, leaving no underlyingly toneless tone-bearing Units that could host a realis-marking H tone.  Only the radical, i.e.\ the first syllable within Gyeli verb stems, is specified for a lexical tone, which is either H or L. Second and third syllables are synchronically and/or diachronically extension morphemes\footnote{Extension morphemes in Bantu languages occur between the verb root and the final vowel and serve valency-changing purposes, encoding categories such as causative, applicative, reciprocal and passive.} and underlyingly toneless. They receive their tonal specification through their phonological environment or the attachment of grammatical tone melodies. In the case of negation with H...-{\itshape lɛ}, the H tone attaches to the left of the verb and spreads rightwards onto the negation suffix -{\itshape lɛ}. Verbs with a H verb radical thus surface all H, as shown in \REF{leH}.\footnote{Monosyllabic H verbs as in \REF{leHL} surface HL, due to a lowering rule for this specific syllable type.}

\begin{exe} \ex\label{leH}
\begin{xlist}
\ex\label{leHL} H $\rightarrow$ H \\
 pɛ̂ `choose' > pɛ́-lɛ́
\ex\label{leHLbi} H {\O} $\rightarrow$ H H \\
 símɛ  `respect' > símɛ́-lɛ́
\ex\label{leHLL} H {\O} {\O} $\rightarrow$ H H H   \\
líyɛlɛ  `show' > líyɛ́lɛ́-lɛ́
\end{xlist}
\end{exe}

Verb radicals that are specified L are featurally changed to H. In monosyllabic verbs, this H tone spreads onto the negation suffix, as shown in \REF{leL}.

\begin{exe}
\ex\label{leL} L $\rightarrow$ H \\
dè `eat' > dé-lɛ́
\end{exe}

In bi- and trisyllabic verb stems with an L radical, however, only the radical is changed to H. In these cases, the H does not spread to the second and third toneless syllables nor to the negation suffix, as illustrated in \REF{leL2}.

\begin{exe} \ex\label{leL2}
\begin{xlist}
\ex\label{leLL} L {\O} $\rightarrow$ H L \\
gyàga `buy' > gyágà-lɛ̀
\ex\label{leLLL} L {\O} {\O} $\rightarrow$ H L L \\
vìdɛga  `turn' > vídɛ̀gà-lɛ̀
\end{xlist}
\end{exe}

The realis/irrealis distinction is not generally lost under negation (see \tabref{tab:3:moods} for the distribution of mood categories). In the case of the past negation markers {\itshape  sàlɛ́/pálɛ́}, the final H tone marks realis, as shown in \REF{meta2b}. In contrast, the future negation marker {\itshape kálɛ̀} with its final L tone belongs in the irrealis category, just like its affirmative counterpart, as illustrated in \REF{SN3xb}. The present is the only tense-mood category that switches from realis in the affirmative to irrealis under negation.  \citet{Miestamo2007} refers to this type as ``A/NonReal''. 

The second paradigmatic asymmetry concerns the neutralization of all (grammaticalized) aspect categories under negation (except for the non-complete accomplishment (\textsc{nca}) {\itshape sílɛ} `finish'). This also includes the inchoative, which formally clusters with tense-mood categories and which is subsumed under the present tense-mood category when negated, as shown in \REF{INCH}.

\begin{exe} 
\ex\label{INCH}
\begin{xlist}
\ex\label{INCHa} 
  \glll     yàá  dé mántúà \\
          yàá dè-H {H\backslash{}ma-ntúà}\\
              1{\PL}.{\INCH} eat-{\R} {{\OBJ}.{\LINK}\backslash{}6-mango}\\
    \trans `We are at the beginning of eating mangoes.'
\ex\label{INCHb}
  \glll     yá/*yàá délɛ́ mántúà \\
            ya-H dè-lɛ {H\backslash{}ma-ntúà}\\
              1{\PL}-{\PRS}/*1{\PL}.{\INCH} eat-{\NEG} {{\OBJ}.{\LINK}\backslash{}6-mango}\\
    \trans `We do not eat mangoes.'
\end{xlist}
\end{exe}

Aspect marking in Gyeli is optional and mostly achieved through auxiliary verbs in affirmative constructions. As \tabref{tab:4:aspect} shows, most aspect markers are restricted to a specific tense-mood category. As a result, the form of the {\STAMP} marker is restricted as well, as generally illustrated by the agreement class 2 {\STAMP} marker. For instance, aspect markers that can only occur in the present tense have always an L tone {\STAMP}. Those restricted to past categories allow the {\STAMP} markers for both recent and remote past, illustrated by  {\itshape bà} and {\itshape báà}. The prospective marker {\itshape múà} is special in that the {\STAMP} clitic for the first and second person singular and agreement class 1 take an L tone, as shown with {\itshape mɛ̀} for the first person singular in \tabref{tab:4:aspect}, while all others have an H tone, shown by {\itshape bá}. The aspect categories that have no tense-mood restriction allow all {\STAMP} forms. Aspect markers without translations in \tabref{tab:4:aspect} have no lexical meaning. 

\begin{table}
\caption{Aspect markers in Gyeli}
\label{tab:4:aspect}
\begin{tabular}{p{3cm}llll}
  \lsptoprule
Status & Aspect  & {\STAMP} &  {\TM}         & Function \\
         & marker   &            & restriction &            \\
\midrule
Grammaticalized  & \itshape nzíí & \itshape bà &  \textsc{prs} & \textsc{prog} \\ 
 verbs & \itshape nzí & \itshape bà, báà &  \textsc{pst} & \textsc{prog} \\
 & \itshape nzɛ́ɛ́ & \itshape bà & \textsc{sub} & \textsc{prog} \\ 
 & {\itshape pã́ } `do first' & all forms &  none & \textsc{prior} \\
 \midrule
Verbs with & {\itshape lɔ́} `come' & \itshape bá &  \textsc{prs} & \textsc{retro} \\ 
lexical meaning & {\itshape bwàá} `have' & \itshape bà, báà &  \textsc{pst} & \textsc{prf}\\
& {\itshape múà} `be' & \itshape mɛ̀, bá &  \textsc{fut} & \textsc{prosp}   \\
  & {\itshape sílɛ}  `finish' & all forms &  none & \textsc{nca} \\
\midrule
Reduplication & stem-stem & \itshape bá &  \textsc{prs} & \textsc{hab} \\
\midrule
Postverbal & \itshape mɔ̀ & \itshape bà &  \textsc{rec.pst} & \textsc{compl} \\
  \lspbottomrule
\end{tabular}
\end{table}

As a consequence of aspect neutralization under negation, the negated sentence is only marked for tense-mood, but not for aspect. The past progressive auxiliary {\itshape nzí}, for instance, is lost under negation and replaced by the past negation auxiliary, as in \REF{PROG}.

\begin{exe} 
\ex\label{PROG}
\begin{xlist}
\ex\label{PROGa} 
  \glll     yà {\bfseries nzí}  dè mántúà \\
          ya nzí dè {H\backslash{}ma-ntúà}\\
              1{\PL}.{\PST} {\PROG}.{\PST} eat {{\OBJ}.{\LINK}\backslash{}6-mango}\\
    \trans `We were eating mangoes.'
\ex\label{PROGb}
  \glll     yà {\bfseries sàlɛ́/pálɛ́} dè mántúà  \\
            ya.{\PST} sàlɛ́/pálɛ́ dè {H\backslash{}-ntúà}\\
              1{\PL}.{\PST} {\NEG}.{\PST} eat {{\OBJ}.{\LINK}\backslash{}6-mango}\\
    \trans `We did not eat mangoes.'
\end{xlist}
\end{exe}

It is impossible to combine aspect and negation markers in a main clause for the reason that more than one grammaticalized auxiliary in a construction, not counting modal semi-auxiliaries,  is prohibited \REF{PROGc}. Also with aspect auxiliaries that are restricted to the present tense-mood category, the attachment of the present negation circumfix H...-{\itshape lɛ}, as in \REF{PROGd}, is not allowed. 

\begin{exe} 
\ex\label{PROG2}
\begin{xlist}
\ex[*]{\label{PROGc}
  \glll     yà sàlɛ́/pálɛ́ nzí/ì dè mántúà\\
            ya sàlɛ́/pálɛ́ nzí/ì dè {H\backslash{}ma-ntúà}\\
              1{\PL}.{\PST} {\NEG}.{\PST} {\PROG}.{\PST} eat {{\OBJ}.{\LINK}\backslash{}6-mango}\\
    \trans `We were not eating mangoes.'}
\ex[*]{\label{PROGd}
  \glll     yà nzílɛ́ dè mántúà\\
            ya nzí-lɛ dè {H\backslash{}ma-ntúà}\\
              1{\PL}.{\PST} {\PROG}.{\PST}-{\NEG} eat {{\OBJ}.{\LINK}\backslash{}6-mango}\\
    \trans `We were not eating mangoes.'}
\end{xlist}
\end{exe}

Aspect and negation can only be combined through embedded clauses with negation expressed in the matrix and aspect expressed in the dependent clause, as shown in \REF{lea1}.

\begin{exe} \ex\label{lea1}
  \glll  mɛ̀ sàlɛ́/pálɛ́ bɛ̀ {\ob}mɛ̀ nzɛ́ɛ́ dè mántúà{\cb} \\
          mɛ sàlɛ́/pálɛ́ bɛ̀ {\db}mɛ nzɛ́ɛ́ dè {H\backslash{}ma-ntúà}\\
          1{\SG}.{\PST} {\PST}.{\NEG} be 1{\SG}.{\PST} {\PROG}.{\SUB} eat {{\OBJ}.{\LINK}\backslash{}6-mango}\\
    \trans `I was not eating mangoes.' (lit. `I was not that I am eating mangoes.')
\end{exe}

\subsection{Negation in non-declaratives}
\label{sec:gye:SNnondecl}

Non-declarative sentences comprise questions as well as imperatives and subjunctives. In Gyeli, questions (\sectref{sec:gye:Q}) follow the standard negation patterns, whereas negated imperatives (\sectref{sec:gye:ti}) and negated subjunctives (\sectref{sec:gye:duu}) require non-standard negation strategies. 

Imperative and subjunctive constructions are generally expected to behave differently under negation. \citet{Nurse2008} posits that Bantu languages tend to have different negation strategies based on sentence types, commonly opposing negation in main versus dependent and indicative versus subjunctive clauses. While Gyeli imperatives and subjunctives indeed require a non-standard negation strategy, one cannot, however, make a case for a neat dichotomy between imperative/subjunctive and indicative clauses. The reason for this is that both the imperative negation with {\itshape tí} and the subjunctive negation with {\itshape dúù} are also found in (some) negated main clauses in the present tense-mood category. Both non-standard negation markers  are asymmetric, including shifts in finiteness and tonal changes. 

\subsubsection{Negation in questions}
\label{sec:gye:Q}

Just like in declarative sentences, the form of the negation marker in questions is determined by the tense-mood category with which it appears (\tabref{tab:2:negmarkers1}). Negated questions are also subject to the same constructional and paradigmatic asymmetries described in \sectref{sec:gye:CA} and \sectref{sec:gye:PA}.  

The type of question, e.g. polar questions or constituent questions, has no impact on the type of negation strategy. Example \REF{PQ1} shows an example of a polar question, where only the question marker {\itshape nà} indicates that it is not an affirmative sentence.   

\ea\label{PQ1}
\ea\label{PQ1a}
 \glll     nà wɛ́ nyɛ́ nyɛ̂  \\
            nà wɛ-H nyɛ̂-H nyɛ̂ \\
             {\Q} 2{\SG}-{\PRS} see-{\R} 1.{\OBJ}  \\
    \trans `Do you see him/her?'
\ex\label{PQ1b}
\glll    nà wɛ̀ɛ́ nyɛ́lɛ́ nyɛ̂  \\
            nà wɛ̀ɛ́ nyɛ̂-lɛ nyɛ̂ \\
             {\Q} 2{\SG}.{\PRS}.{\NEG} see-{\NEG} 1.{\OBJ}  \\
    \trans `Don't you see him/her?'
\z
\z

Constituent questions employ animacy-based interrogative pronouns for all argument types, as for instance an object in a double object construction in \REF{QIO}. The negated counterpart follows the standard negation strategy (\sectref{sec:gye:SN}) with different negation markers for different tense-mood categories and both constructional and paradigmatic asymmetries.  

\ea\label{QIO}
\ea\label{QIOa}
  \glll  nzá àà vɛ̂ béfùmbí \\
         nzá àà vɛ̂ {H\backslash{}be-fùmbí}\\
         who 1.{\FUT} give {{\OBJ}.{\LINK}\backslash{}8-orange}\\
    \trans `To whom will s/he give oranges?'
\ex\label{QIOb}
  \glll  nzá àà kálɛ̀ vɛ̂ béfùmbí \\
         nzá àà kálɛ̀ vɛ̂ {H\backslash{}be-fùmbí}\\
         who 1.{\FUT} {\NEG}.{\FUT} give {{\OBJ}.{\LINK}\backslash{}8-orange}\\
    \trans `To whom won't s/he give oranges?'
\z    
\z

\subsubsection{Non-standard negation with {\itshape tí}}
\label{sec:gye:ti}

The primary use of {\itshape tí} is to negate imperatives.\footnote{I could not discern any clear lexical meaning for {\itshape tí}. Out of context, speakers will translate it as `without', but this is probably because the derived use of {\itshape tí} as a privative preposition, as discussed below, is the only case where a lexical translation equivalent can be found.} Positive imperatives have three different construction types, as shown in \REF{IMPpos}. The verbal imperative form is always marked by an HL tone pattern (\tabref{Tab:TM-Tone}). In \REF{IMPposa}, the addressee is a single individual and the imperative verb occurs without any person marking. In contrast, if there are multiple addressees, a postverbal plural particle is used, as in \REF{IMPposb}. In a cohortative construction, as in \REF{IMPposc}, both the first person plural subject marker and the postverbal plural particle are used.

\ea\label{IMPpos} 
    \ea\label{IMPposa} {\bfseries dê} \\
        \trans `Eat (\textsc{sg})!' 
    \ex\label{IMPposb} {\bfseries dê} ŋgà \\
        \trans `Eat (\textsc{pl})!'
    \ex\label{IMPposc} 
        \glll yá {\bfseries dê} ŋgà \\
        ya-H dê ŋga \\
        1{\PL}-{\PRS} eat.{\SBJV} {\PL} \\
        \trans `Let's eat!'
    \z
\z

In negated imperatives, as in \REF{IMPneg}, the lexical verb is preceded by the negative auxiliary {\itshape tí}. Negated imperatives are subject to a constructional asymmetry pertaining to finiteness that is shifted from the lexical verb to the auxiliary. This shift in finiteness formally manifests in two ways. First, the HL pattern on positive imperatives is lost and the lexical verb surfaces with its underlying lexical tones. Second, the plural particle follows the auxiliary and not the lexical verb. 

\ea\label{IMPneg} 
    \ea\label{IMPnega} {\bfseries tí} {\bfseries dè} \\
        \trans `Don't (\textsc{sg}) eat!'
    \ex\label{IMPnegb} {\bfseries tí} ŋgá {\bfseries dè} \\
        \trans `Don't (\textsc{pl}) eat!'
    \ex\label{IMPnegc} 
        \glll yá {\bfseries tí} ŋgá {\bfseries dè} \\
              ya-H tí ŋga dè \\
        	1{\PL}-{\PRS} {\NEG} {\PL} eat \\
        \trans `Let's not eat!'
    \z
\z

The negative auxiliary {\itshape tí} is also used in three non-imperative constructions.  First, it appears as a negation marker in present tense main clauses.
The choice between the standard present tense negation circumfix H...-{\itshape lɛ} (\sectref{sec:gye:SN}) and the non-standard negation auxiliary {\itshape tí} in present tense main clauses relates to information structure principles and an immediate-after-verb focus position.\footnote{There are two ways to classify {\itshape tí} with respect to its status as standard or non-standard negation marker. Since {\itshape tí} is a productive means for negating declarative verbal main clauses (in the present tense), it could be analyzed as standard negation. However,
since it is part of a special focus construction, it is not the neutral way of negating the present tense, which speaks against treating it as standard negation.
I adopt the latter analysis.} If the standard negation circumfix is used, as in \REF{tilea}, the negation itself, thus the truth value, is in focus. This could be, for instance, as a contradiction to a statement such as `You are eating'. In contrast, if {\itshape tí} is used, as in \REF{tileb}, the lexical verb is in focus position. The two constructions therefore serve to disambiguate the scope of negation; more information on scope under negation is provided in \sectref{sec:gye:Scope}. 

\ea\label{tile}
\ea\label{tilea}
  \glll  mɛ̀ɛ́ dé{\bfseries lɛ́} \\
        mɛ̀ɛ́ dé-lɛ́ \\
         1{\SG}.{\PRS}.{\NEG} eat-{\NEG}    \\
    \trans `[You are eating.] -- I \textbf{don't} eat.'  
\ex\label{tileb}
  \glll  mɛ̀ {\bfseries tí} dè\\
         mɛ tí dè \\
          1{\SG} {\NEG}  eat   \\
    \trans `I don't \textbf{eat} [I drink].' 
\z
\z

Second, {\itshape tí} is used in nominal clauses as the privative preposition `without' (\sectref{sec:gye:Deriv}), and third, {\itshape tí} serves as a negative  subordinator (\sectref{sec:gye:Coord}). 

\subsubsection{Non-standard negation with {\itshape dúù}}
\label{sec:gye:duu}

The second non-standard negation marker {\itshape dúù} is a true auxiliary in Gyeli, which is defined by its tense-mood restriction to the present and subjunctive category. This is different from other modal verbs, such as {\itshape wúmbɛ} `want' or {\itshape lèmbo} `know', which can occur in every tense-mood category and are thus semi-auxiliaries. The marker {\itshape dúù} serves as the negative counterpart of two affirmative constructions.  First, it negates the modal verb {\itshape yánɛ} `must' by replacing it, as shown in \REF{duu1}. This construction shows constructional asymmetry in that the modal of the affirmative is replaced by the negative auxiliary. The positive and the negative are aligned, however, in the mood category, as both take the realis-marking H tone. 

\ea\label{duu1}
\ea\label{duu1a}
  \glll bé dúú vũ̀ũ̀\\
      be-H dúù-H vũ̀ũ̀ \\
        2{\PL}-{\PRS} must.not-{\R} worry \\
    \trans `You (\textsc{pl}) should/must not worry.'
\ex\label{duu1b}
  \glll bé yánɛ́ vũ̀ũ̀ \\
      be-H yánɛ-H vũ̀ũ̀ \\
        2{\PL}-{\PRS} must-{\R} worry \\
    \trans `You (\textsc{pl}) should/must worry.'
\z
\z

Second, the auxiliary {\itshape dúù}  is also  used to negate subjunctive clauses, as in \REF{duu2}. These subjunctive clauses are constructionally asymmetric, since the finiteness is shifted from the lexical verb in the affirmative to the auxiliary under negation. 

\ea\label{duu2}
\ea\label{duu2a}
  \glll bé dúù kɛ̀ tísɔ̀nì \\
      be-H dúù kɛ̀ tísɔ̀nì \\
        2{\PL}-{\PRS} must.not.{\SBJV} go $\emptyset$.7.town \\
    \trans `You (\textsc{pl}) may/should not go to town.'
\ex\label{duu2b}
  \glll bé kɛ́ɛ̀ tísɔ̀nì \\
      be-H kɛ́ɛ̀ tísɔ̀nì \\
        2{\PL}-{\PRS} go.{\SBJV} $\emptyset$.7.town \\
    \trans `You (\textsc{pl}) may/should go to town.'
\z
\z

The use of {\itshape dúù} as negating a  modal or the subjunctive  is clearly marked by its distinct tonal patterns. In modal negation in \REF{duu1a}, it takes the realis-marking H tone of the present tense-mood category. In  subjunctive negation, as in \REF{duu2a}, it surfaces with the subjunctive HL pattern without the realis-marking H tone. 

Just like the other non-standard negation marker {\itshape tí}, {\itshape dúù} is also used in dependent clauses (\sectref{sec:gye:non-main}). 

\subsection{Negation in stative predications}
\label{sec:gye:Stative}

Gyeli has four types of non-verbal predications in affirmative present tense clauses, following \posscitet{Payne1997} classification of non-verbal clauses. In Gyeli, the  different types are marked by different copulas and/or differences in constituent order. In other tenses than the present and under negation, these copulas are replaced by an inflected form of the verb {\itshape bɛ̀}  `be'. There is no difference in the encoding of permanent versus temporary property assignment, neither in affirmative nor negated sentences. 

The most frequent copula in present affirmative non-verbal clauses is the ``{\STAMP} copula'', which has the same segmental form as long-vowel {\STAMP} clitics (\tabref{Tab:TM-Tone}), marking subject agreement. The {\STAMP} copulas for the first and second person singular and agreement class 1 have an H tone pattern, while all others have an HL pattern. As shown in \REF{COP}, the {\STAMP} copula marks relations of equation, proper inclusion, attribution, and locative predication between the subject and the predicate. As such, the predicates of a {\STAMP} copula construction include nouns, adjectives, and locative noun phrases. 
(In other tenses than the present, the specific tense form of the verb {\itshape bɛ̀} `be' is used for all of these constructions.)


\ea\label{COP}
\ea{\label{COPa} Equation\\
  \gll    Àdà {\bfseries àà} sã̂ wã̂   \\ 
          $\emptyset$.1.{\PN} 1.{\COP} $\emptyset$.1.father 1-{\POSS}.1\textsc{sg}\\}
    \trans `Ada is my father.'
\ex{\label{COPb} Proper inclusion\\
  \gll    Àdà {\bfseries àà} ŋgɛ̀lɛ́nɛ̀   \\ 
          $\emptyset$.1.{\PN} 1.{\COP} $\emptyset$.1.teacher \\}
    \trans `Ada is a teacher.'
\ex{\label{COPc} Attribution\\
  \gll    Àdà {\bfseries àà} mpà   \\ 
          $\emptyset$.1.{\PN} 1.{\COP} good \\}
    \trans `Ada is good.'
\ex{\label{COPd} Locative predication\\
  \gll    Àdà {\bfseries àà} ndáwɔ̀ dé tù  \\ 
          $\emptyset$.1.{\PN} 1.{\COP} $\emptyset$.9.house {\LOC} inside \\}
    \trans `Ada is in the house.'
    \z 
\z 

Negation of these non-verbal predications patterns with standard negation strategies using the verb {\itshape bɛ̀}  `be', as shown in \REF{COPx}.

\ea\label{COPx}
\ea{\label{COPxa} Equation \\
  \gll    Àdà {\bfseries àá} {\bfseries bɛ́-lɛ́}  sã̂ wã̂   \\
          $\emptyset$.1.{\PN} 1.{\PRS}.{\NEG} be-{\NEG} $\emptyset$.1.father 1-{\POSS}.1\textsc{sg}\\}
    \trans `Ada is not my father.'
\ex{\label{COPxb} Proper inclusion\\
  \gll    Àdà {\bfseries àá} {\bfseries bɛ́-lɛ́} ŋgɛ̀lɛ́nɛ̀\\
          $\emptyset$.1.{\PN} 1.{\PRS}.{\NEG} be-{\NEG} $\emptyset$.1.teacher\\}
    \trans `Ada is not a teacher.'
\ex{\label{COPxc} Attribution\\
  \gll    Àdà {\bfseries àá} {\bfseries bɛ́-lɛ́} mpà\\
          $\emptyset$.1.{\PN} 1.{\PRS}.{\NEG} be-{\NEG} good\\}
    \trans `Ada is not good.'
\ex{\label{COPxd} Locative predication\\
  \gll    Àdà {\bfseries àá} {\bfseries bɛ́-lɛ́} ndáwɔ̀ dé tù\\
          $\emptyset$.1.{\PN} 1.{\PRS}.{\NEG} be-{\NEG} $\emptyset$.9.house {\LOC} inside\\}
    \trans `Ada is not in the house.'
\z
\z

Existential clauses are the second type of non-verbal predicates. While they also use the {\STAMP} copula, their predicates are typically adverbs, as in \REF{EXa}. Under negation, the adverb is dropped. In contrast to the constructions in \REF{COPx}, negated existentials allow for the subject noun to appear either clause-finally, as in \REF{EXb}, or clause-initially, as in \REF{EXc}. Example \REF{EXb} represents the unmarked order as either a thetic sentence or with the postverbal noun being in the dedicated focus position. If the negated verb appears phrase-finally, the focus is on the negation and thus on the truth value of the clause. 

\ea\label{EX}
\ea\label{EXa} 
  \gll    b-ùdã̂ báà tè\\
          ba.2-woman 2.{\COP} there \\
    \trans `There are women.'
\ex\label{EXb} 
  \gll    bá bɛ́-lɛ́ b-ùdã̂ \\
          2.{\PRS} be-not 2-woman \\
    \trans `There are no \textbf{women} (or thetic).'
\ex\label{EXc} 
  \gll    b-ùdã̂ bá bɛ́-lɛ́\\
          2-woman 2.{\PRS} be-{\NEG} \\
    \trans `There \textbf{are no} women.'
\z
\z

Third, Gyeli marks identificational clauses with the identificational marker {\itshape wɛ́}. In contrast to the {\STAMP} copula, it is invariable and does not agree with the subject. The identificational marker typically links a nominal subject to a demonstrative, as in \REF{we1a}, or an anaphoric marker, as in \REF{we1b}, which is uttered at the end of a storytelling. 

\ea\label{we1}
\ea\label{we1a}
  \gll ntɛ́mbɔ́ wã̂ {\bfseries wɛ́} nû \\
    $\emptyset$.1.younger.opposite.sex.sibling 1-{\POSS}.1\textsc{sg} {\ID} 1.{\DEM}.{\PROX}  \\
    \trans `This is my younger brother/sister.'
\ex\label{we1b}
  \gll kàndá {\bfseries wɛ́} ndɛ̀ \\
        $\emptyset$.7.proverb {\ID} {\ANA}  \\
    \trans `This is the story.'  
\z
\z

The negated counterparts of identificational constructions are also encoded by the verb {\itshape bɛ̀} `be', as shown in \REF{we2}. As with negated existentials, negated identification clauses with a demonstrative allow for more freedom in constituent order, changing the topic referent. In \REF{we2a}, the discourse topic is the sibling, while the deictic referent of the demonstrative is in focus position. In \REF{we2b}, on the other hand, the discourse topic is the deictic referent of the demonstrative, while the nominal complement {\itshape ntɛ́mbɔ́ wã̂} `my younger brother/sister' constitutes new information.\footnote{The demonstrative in \REF{we2b} replaces the {\STAMP} clitic segmentally, but takes the tonal pattern that the {\STAMP} clitic would take in this environment.}

\ea\label{we2}
\ea\label{we2a}
  \gll ntɛ́mbɔ́ wã̂ {\bfseries àá} {\bfseries bɛ́-lɛ́} nû \\
          $\emptyset$.1.younger.opposite.sex.sibling 1-{\POSS}.1\textsc{sg} 1.{\PRS}.{\NEG} be-{\NEG} 1.{\DEM}.{\PROX}  \\
    \trans `My younger brother/sister is not this one.'
\ex\label{we2b}
  \gll nú bɛ́-lɛ́ ntɛ́mbɔ́ wã̂\\
  1.{\DEM}.{\PRS} be-{\NEG} $\emptyset$.1.younger.opposite.sex.sibling 1-{\POSS}.1\textsc{sg}\\
      \trans `This is not my younger brother/sister.'
\ex\label{we2c}
  \gll kàndá {\bfseries yí} {\bfseries bé-lɛ́} {\bfseries yí}-ndɛ̀ \\
        $\emptyset$.7.proverb 7.{\PRS} be-{\NEG}  7-{\ANA}  \\
    \trans `This is not the story.'  
\z
\z

While the anaphoric marker (which is distinct from the proximal and distal demonstratives) can appear bare with the identificational marker {\itshape wɛ́} in \REF{we1b}, it requires an agreement prefix in the negated construction in \REF{we2c}. 

Finally, possessive predication is expressed by the combination of a verb and comitative marker {\itshape bɛ̀ nà}  `be with'.  As such, these constructions generally follow the same pattern as negation in main clauses, as shown in \REF{honey}. At the same time, they allow a choice between the standard negation form in \REF{honey2} and the use of the non-standard privative marker {\itshape tí}  (\sectref{sec:gye:ti} and \sectref{sec:gye:Deriv}), as in \REF{honey3}, which seems to emphasize the absence of something. 

\ea \label{honey}
\ea \label{honey1}
  \glll  mɛ́ {\bfseries bɛ́} {\bfseries nà} nkwànò  \\
        mɛ-H bɛ̀-H nà nkwànò \\
          1{\SG}-\textsc{prs} be-{\R} {\COM} $\emptyset$.3.honey\\
    \trans `I have honey.'
\ex \label{honey2}
  \glll  mɛ́ {\bfseries bɛ́lɛ́} {\bfseries nà} nkwànò  \\
        mɛ-H bɛ̀-lɛ nà nkwànò \\
          1{\SG}-\textsc{prs} be-{\NEG} {\COM} $\emptyset$.3.honey\\
    \trans `I don't have any honey.'
\ex \label{honey3}
  \glll  mɛ̀ {\bfseries tí} {\bfseries bɛ̀} {\bfseries nà} nkwànò  \\
        mɛ  tí bɛ̀ nà nkwànò \\
          1{\SG}.{\SBJ} {\NEG} be {\COM} $\emptyset$.3.honey\\
    \trans `I'm without honey.'    
\z
\z

\subsection{Negation in non-main clauses}
\label{sec:gye:non-main}

Gyeli does not have any dedicated negation strategies that make distinctions between main and non-main clauses. Just like in main clauses (\sectref{sec:gye:SN}), non-main clauses allow for standard and non-standard negation constructions.
Thus the dichotomy of different negation types contrasting indicative/subjunctive or main/relative clauses that has been observed in many Bantu languages \citep{Nurse2008} does not apply in Gyeli. Example \REF{SREL} shows standard negation in a relative clause with the same constructional and paradigmatic asymmetries described in \sectref{sec:gye:CA}
 and \sectref{sec:gye:PA}.

\ea\label{SREL} 
\ea\label{SREL1}
  \glll bá dyúwɔ́ lɛ́kɛ́lɛ̀ {\ob}lé {\bfseries wɛ́} {\bfseries làwɔ̀}{\cb}\textsubscript{{\REL}}\\
      ba-H dyúwɔ-H {H\backslash{}lɛ-kɛ́lɛ̀} {\db}lé wɛ-H làwɔ\\
         2-{\PRS} understand {{\OBJ}.{\LINK}\backslash{}5-language} {\db}5:{\ATTR} 2{\SG}-{\PRS} speak\\
    \trans `They understand the language that you speak.'
\ex\label{SREL2} 
  \glll bá dyúwɔ́ lɛ́kɛ́lɛ̀ {\ob}lé {\bfseries wɛ̀ɛ́} {\bfseries láwɔ̀lɛ̀}{\cb}\textsubscript{{\REL}}\\
        ba-H dyúwɔ-H {H\backslash{}lɛ-kɛ́lɛ̀} {\db}lé wɛ̀ɛ́ làwɔ-lɛ\\
         2-{\PRS} understand {{\OBJ}.{\LINK}\backslash{}5-language} {\db}5:{\ATTR} 2{\SG}.{\PRS}.{\NEG} speak-{\NEG}\\
    \trans `They understand the language that you don't speak.'
\z
\z

Standard negation is also used in complement clauses with say- or think-verbs, as in \REF{COMP3} and its negative counterpart in \REF{COMP3x}. 

\ea\label{COMP3} 
  \glll ndí wɛ́ lèmbó {\ob}nâ mbvúndá {\bfseries nyíì} {\bfseries bvúdà} nà mbvúndá{\cb}\textsubscript{{\COMP}}\\
        ndí wɛ-H lèmbo-H {\db}nâ mbvúndá nyíì bvúda nà mbvúndá\\
         but 2{\SG}-{\PRS} know-{\R} {\db}{\COMP} $\emptyset$.9.trouble 9.{\FUT} fight {\COM} $\emptyset$.9.trouble\\
    \trans `But you know that violence will create more violence.'
\z

\ea\label{COMP3x} 
  \glll ndí wɛ́ lèmbó {\ob}nâ mbvúndá {\bfseries nyíì} {\bfseries kálɛ̀} {\bfseries bvúdà} nà mbvúndá{\cb}\textsubscript{{\COMP}}\\
        ndí wɛ-H lèmbo-H {\db}nâ mbvúndá nyíì kálɛ̀ bvúda nà mbvúndá\\
         but 2{\SG}-{\PRS} know-{\R} {\db}{\COMP} $\emptyset$.9.trouble 9.{\FUT} {\NEG}.{\FUT} fight {\COM} $\emptyset$.9.trouble\\
    \trans `But you know that violence will not create more violence.'
\z

There are non-main clauses with the non-standard negation markers {\itshape dúù} and {\itshape tí} to mark predication negation. The latter is special in that it can also act as the negative subordinator of an infinitival clause (\sectref{sec:gye:Coord}).  If a complement clause expresses wishes, orders, obligations, or prohibitions, the subjunctive is used, which requires the non-standard negation marker {\itshape dúù} under negation, as in \REF{SBJV}.

\ea\label{SBJV}
\ea\label{SBJV1}
  \glll  yíì mpìndá nâ wɛ́ djíw{\bfseries ɔ́ɔ̀} bésâ \\
          yíì mpìndá nâ wɛ-H djíwɔ́ɔ̀ {H\backslash{}be-sâ}\\
              7.{\COP} $\emptyset$.9.prohibition {\COMP} 2{\SG}-{\PRS} steal.{\SBJV} {{\OBJ}.{\LINK}\backslash{}8-thing}\\
    \trans `It's forbidden that you steal things.'
\ex\label{SBJV2}
  \glll  yíì mpìnàgà nâ wɛ́ {\bfseries dúù} djíwɔ̀ bésâ\\
          yíì mpìnàgà nâ wɛ-H dúù djíwɔ {H\backslash{}be-sâ}\\
              7.{\COP} $\emptyset$.3.obligation {\COMP} 2{\SG}-{\PRS} must.not.{\SBJV} steal {{\OBJ}.{\LINK}\backslash{}8-thing}\\
    \trans `It's necessary that you don't steal things.'    
\z
\z

Furthermore, {\itshape dúù} can negate a modal in dependent clauses, as in \REF{SBJVx}.

\ea\label{SBJVx}
\ea\label{SBJVxa}
  \glll kálɛ̀ mɛ̀ báà kì nâ bá {\bfseries yánɛ́} bɛ̀ bédéwɔ̀ \\
         kálɛ̀ mɛ̀ báà kì nâ ba-H yánɛ-H bɛ̀ H\backslash{}be-déwɔ̀ \\
      {\NEG} 1{\SG}  2.{\FUT} say {\COMP} 2-{\PRS} must-{\R} grow {{\OBJ}.{\LINK}\backslash{}8-thing} \\
    \trans `It's not me, they [= who] will say that they must grow food.'
\ex\label{SBJVxb}
  \glll kálɛ̀ mɛ̀ báà kì nâ bá {\bfseries dúú} bɛ̀ bédéwɔ̀ \\
         kálɛ̀ mɛ̀ báà kì nâ ba-H dúù-H bɛ̀ H\backslash{}be-déwɔ̀ \\
      {\NEG} 1{\SG}  2.{\FUT} say {\COMP} 2-{\PRS} must.not-{\R} grow {{\OBJ}.{\LINK}\backslash{}8-thing} \\
    \trans `It's not me, they [= who] will say that they must not grow food.'
\z
\z

In complement clauses with the subjunctive, negation is usually raised from the dependent to the main clause, as shown in \REF{duunat}. Negative raising is possible and frequent with the subjunctive, since the subjunctive in Gyeli inherently relates to volition, as reflected also in the predicate of the main clause. Volition predicates constitute one class of predicates involved in negative raising \citep{Horn1989}. Negative raising with other types of verbs that do not entail a subjunctive form in the dependent clause, such as for instance say-verbs, is not possible as the meaning would change. 

\begin{exe} \ex\label{duunat}
  \glll  {\bfseries mɛ̀ɛ́} wúmbɛ́{\bfseries lɛ́} nâ wɛ́  djíw{\bfseries ɔ́ɔ̀} bésâ \\
          mɛ̀ɛ́ wúmbɛ-lɛ nâ wɛ-H djíwɔ́ɔ̀ {H\backslash{}be-sâ}\\
             1.{\PRS}.{\NEG} want-{\NEG} {\COMP} 2{\SG}-{\PRS} steal.{\SBJV} {{\OBJ}.{\LINK}\backslash{}8-thing}\\
    \trans `I don't want that you steal things.'
\end{exe}

\subsection{Negative lexicalizations}

Negative lexicalizations in Gyeli are rare. The most common one is the negative modal auxiliary {\itshape dúù} `must not'. Its main use is in negative subjunctives and therefore discussed in \sectref{sec:gye:duu}.
One may argue that the noun {\itshape mpìndá}  `prohibition' in \REF{SBJV1} is another instance of negative lexicalization, as it generates an implicature that the action in the following complement clause is not to be done. The noun does not, however, perform clausal negation in a syntactic sense.

\section{Non-clausal negation}
\label{sec:gye:3}

While Gyeli has several negation markers for clausal negation, non-clausal negation is rather restricted.
After discussing pro-sentence forms of negative replies in \sectref{sec:gye:NR}, I turn to the lack of negative indefinites and quantifiers in \sectref{sec:gye:INDEF} and the absence of negative derivation in \sectref{sec:gye:Deriv}. 

\subsection{Negative replies}
\label{sec:gye:NR}

Gyeli has several pro-sentence forms for agreement and disagreement signal. Example \REF{yes} provides a list of pro-forms that signal agreement and that can be used in an answer to any yes/no-question. They have pragmatic differences. For instance, while \REF{yes1} is the default form, \REF{yes6} is more emphatic. 

\begin{exe} 
\ex\label{yes}
    \begin{xlist}
    \ex\label{yes1} ɛ́ɛ̀ `yes'
    \ex\label{yes2} áà `yes' 
    \ex\label{yes3} èè `yes' 
    \ex\label{yes4} ɛ́ɛ́ `yes' 
    \ex\label{yes5} m̀ḿḿ `yes'
    \ex\label{yes6} ɛ̀hɛ́ɛ́  `yes'
    \end{xlist}
\end{exe}

There are fewer pro-forms for disagreement than for agreement. The default form is {\itshape tɔ̀sâ} in \REF{no1}, which is derived from the negative polarity item {\itshape tɔ̀} and the noun {\itshape sâ} `thing'.

\begin{exe}
\ex\label{no}
    \begin{xlist}
    \ex\label{no1} tɔ̀sâ `no'
    \ex\label{no2} ɛ́'ɛ̂ `no'
    \ex\label{no3} ḿ'm̂ `no'
    \end{xlist}
\end{exe}

Semantically, negative replies match the question's polarity, as shown in (\ref{Qa}--\ref{Qb}).  \citet{EnfieldEtAl2019} refer to this interjection-type  answer as a `yes/no system' (contrasting an `agree/disagree system' in which the answer relates to the truth value of the proposition).

\ea\label{Qa}
    \begin{xlist}[A2: ]
    \exi{Q: }
        \glll {\bfseries nà} wɛ̀ nyɛ́ nyɛ̂ \\
              nà wɛ nyɛ̂-H nyɛ̂ \\
              {\Q}   2\textsc{sg}.{\textsc{rec.}\PST} see-{\R} 1.{\OBJ}\\
        \trans `Did you see him/her?'
    \exi{A1: }\label{yesx}  ɛ́ɛ̀ `yes! [I did see him/her]'
    \exi{A2: }\label{nox} tɔ̀sâ `no! [I didn't see him/her]'
    \end{xlist}
\z 

\ea\label{Qb}
    \begin{xlist}[A2: ]
    \exi{Q: }
        \glll {\bfseries nà} wɛ̀ɛ́ nyɛ́lɛ́ nyɛ̂ \\
              nà wɛ̀ɛ́ nyɛ̂-lɛ nyɛ̂ \\
              {\Q} 2\textsc{sg}.{\PRS}.{\NEG} see-{\NEG} 1.{\OBJ}\\
        \trans `Didn't you see him/her?'
    \exi{A1: }\label{yesx1}  ɛ́ɛ̀ `yes! [I did see him/her]'
    \exi{A2: }\label{nox2} tɔ̀sâ `no! [I didn't see him/her]'
    \end{xlist}
\z 


Interjection-type answers seem more frequent and the default, but repetition-type answers are also used, which echo the verb of the question. The repeated proposition either occurs by itself or in combination with the interjection-type answer, for instance as in {\itshape I didn't see him!} or {\itshape No, I didn't see him}. This combination of a `yes/no' and an `echo system' is the most frequent among mixed answering systems \citep{Moser2018}. The distribution of different answer types has not been investigated at length in Gyeli, but it seems that repetition-type answers and their combinations with an interjection are linked to emphasis and assertion. 

\subsection{Negative indefinites and quantifiers}
\label{sec:gye:INDEF}

Gyeli does not have any negative indefinite pronouns. Equivalents of {\itshape nobody} or {\itshape nothing} in English are expressed by bare nouns with the singular form of {\itshape mùdì} `person' for animates and {\itshape sâ} `thing' for inanimates, as shown in \REF{OBJ} for these bare nouns in object position. 

\ea\label{OBJ}
  \glll   mɛ̀ɛ́ nyɛ́lɛ́ mùdì/sâ \\
         mɛ̀ɛ́ nyɛ̂-lɛ m-ùdì/sâ \\
    1\textsc{sg}.{\PRS}.{\NEG} see-{\NEG} $\emptyset$.1-person/$\emptyset$.7.thing \\
    \trans `I don't see a (the) person/a (the) thing.'
\z    

As the translation suggests, this clause is ambiguous as to the referent of the bare noun. The negative polarity item {\itshape tɔ̀} can optionally be used to clarify when the intended meaning is {\itshape anyone/anything}, as discussed in (\sectref{sec:gye:NegPol}).

If these default bare nouns function as the subject, a negated existential construction (\sectref{sec:gye:Stative}) is used in combination with a relative clause, as in \REF{SBJ}.\footnote{It is cross-linguistically common to use negative existentials to encode negative indefinites \citep{Haspelmath1997,Veselinova2013neg.exist}.}

\ea\label{SBJ}
\ea\label{SBJa}
  \glll bá bɛ́lɛ́ bùdì {\ob}bà kɛ́ tísɔ̀nì{\cb}\textsubscript{{\REL}}\\
  ba-H bɛ̀-lɛ b-ùdì {\db}ba kɛ̀-H tísɔ̀nì\\
  2-{\PRS} be-{\NEG} 2-person {\db}2.{\textsc{rec.}\PST} go-{\R} $\emptyset$.7.town\\
      \trans `Nobody went to town.'
\ex\label{SBJb}
 \glll bà bɛ́lɛ́ bùdì {\ob}bà kɛ́ tísɔ̀nì{\cb}\textsubscript{{\REL}}\\
  ba bɛ̀-lɛ b-ùdì {\db}ba kɛ̀-H tísɔ̀nì\\
  2.{\textsc{rec.}\PST} be-{\NEG} 2-person {\db}2.{\textsc{rec.}\PST} go-{\R} $\emptyset$.7.town\\
      \trans `There \textbf{were} no people who went to town.'
\z
\z

In contrast to the object position, the plural form is preferred with the subject bare noun, although the singular is also possible. In terms of tense marking, the negative clause could also appear in other tense-mood categories, for instance in the past or future. However, the present marking in the negated main clause and potentially a different tense marking in the relative clause is the default. If the main clause was marked for a different tense category, that tense category would be in focus, as in \REF{SBJb} (while the relative clause keeps the past tense marking as well). 

There are also no negative quantifiers such as {\itshape no} or {\itshape none}.  Equivalents of these are either expressed by negated existential clauses, as in \REF{SBJ}, or with {\itshape tí} constructions, as in \REF{NEGwo1b}.  

\subsection{Negative derivation}
\label{sec:gye:Deriv}

Gyeli lacks the diversity in negation (derivational) affixes found in many Indo-European languages that are associated with semantic categories such as ``contradictory negation'' ({\itshape non}-), ``contrary negation'' and ``reversal'' ({\itshape un}-), ``privative negation'' and ``removal'' ({\itshape dis}-), and ``opposition'' ({\itshape anti}-) \citep{CartoniLefer2011,Joshi2020}. 
There are only two instances of negative derivation in Gyeli. One is the negative reply {\itshape tɔ̀sâ} `no', which is a lexicalized form derived from the negative polarity item {\itshape tɔ̀} (\sectref{sec:gye:NegPol}) and the noun {\itshape sâ} `thing'.

The other is the use of the non-standard negation marker {\itshape tí}  with nominal heads. While {\itshape tí} is primarily used as a negation marker for imperatives (\sectref{sec:gye:ti}), it has also grammaticalized into a nominal negation marker that drops the verbal elements of an infinitival clause {\itshape tí bɛ̀ nà} `not being with' to serve as a privative preposition {\itshape tí} `without', as shown in \REF{NEGwo1b} -- the elements in parentheses are usually omitted.

\ea\label{NEGwo1} 
\ea\label{NEGwo1a} 
  \glll  mɛ́ nyúlɛ́ kɔ̀fí {\bfseries nà} ngùɔ́ \\
            mɛ-H nyúlɛ-H kɔ̀fí nà ngùɔ́ \\
         1{\SG}-{\PRS} drink-{\R} $\emptyset$.7.coffee {\COM} $\emptyset$.7.sugar \\
    \trans `I drink coffee with sugar.'
\ex\label{NEGwo1b} 
  \glll  mɛ́ nyúlɛ́ kɔ̀fí {\bfseries tí} ({\bfseries bɛ̀} {\bfseries nà}) ngùɔ́ \\
            mɛ-H nyúlɛ-H kɔ̀fí tí bɛ̀ nà ngùɔ́ \\
         1{\SG}-{\PRS} drink-{\R} $\emptyset$.7.coffee {\NEG} be {\COM} $\emptyset$.7.sugar \\
    \trans `I drink coffee without sugar.'
\z
\z

There is one unique case of reversal opposition between {\itshape jì} `open' and {\itshape jì-bɔ} `close' (lit. `un-open'). The extension suffix -{\itshape bɔ} and its alternate form -{\itshape wɔ} cannot, however, be associated with a general reversal meaning \citep[279]{Grimm2021}.

As Gyeli only has a very small class of adjectives, morphologically negated adjectives do not exist; they would be expressed as other non-verbal predicates (\sectref{sec:gye:Stative}).
Further, Gyeli does not have any case marking so that negation phenomena in connection with case do not apply. 

\section{Other aspects of negation}
\label{sec:gye:4}

\subsection{The scope of negation}
\label{sec:gye:Scope}

Negation can have scope over all constituents in a sentence. Therefore, a sentence as in \REF{Negscope} is ambiguous with respect to which constituent is being negated.

\ea\label{Negscope}
  \glll   Àdà àá gyámbɔ́lɛ́ békwàndɔ̀ kísínì dé tù\\
     Àdà àá gyámbɔ-lɛ {H\backslash{}be-kwàndɔ̀} kísínì dé tù\\
    $\emptyset$.1.{\PN} 1.{\PRS}.{\NEG} cook-{\NEG} {{\OBJ}.{\LINK}\backslash{}8-plantain} $\emptyset$.1.kitchen {\LOC} inside \\
    \trans `Ada doesn't cook plantains in the kitchen.'
\z 

The negation can scope over any {\NP}: the subject (it isn't Ada, it is Mambi), the object (it isn't plantains, it is potatoes), and the adjunct (it isn't in the kitchen, it is outside). It can also negate both the verb (he doesn't cook the plantains, he grills them) and the TAM category (he doesn't cook plantains right now, he already cooked them). 

There are means to disambiguate the focus of negation, following the language's information packaging principles. The most important means of conveying information structure in Gyeli is word order. Prosody, such as stress, does not play a role. Instead, Gyeli has dedicated information structure positions \citep{GuldemannEtAl2015},  with the subject as the default topic and the object as default focus \citep{Grimmforth}. 
Predicate focus is achieved through (pronominal) object fronting as in \REF{PCF}.  

\ea\label{PCF}
\glll à pálɛ́ byɔ̂ gyámbɔ̀\\
    a pálɛ́ byɔ̂ gyámbɔ\\
    1.{\textsc{rec.}\PST} {\PST}.{\NEG} 8.{\OBJ} cook\\
\trans `S/he did not \textbf{cook} them [the plantains].'
\z

Subject and adjunct focus are expressed by cleft constructions, as exemplified for subject focus in \REF{cleftFOC}.

\ea\label{cleftFOC}
\glll yíì pálɛ́ Àdà à gyámbɔ́ békwàndɔ̀\\
  yíì pálɛ́ Àdà a gyámbɔ-H {H\backslash{}be-kwàndɔ̀}\\
  7.{\COP} {\PST}.{\NEG} $\emptyset$.1.{\PN}   1.{\textsc{rec.}\PST} cook-{\R} {{\OBJ}.{\LINK}\backslash{}8-plantain}\\
\trans `It wasn't Ada who cooked the plantains.'
\z

A special focus construction using different negation strategies was discussed in \sectref{sec:gye:ti} in example \REF{tile}. 

Negation also has scope effects in terms of narrow and wide scope readings in negated sentences (see also \sectref{sec:gye:INDEF}).  In \REF{NPI1}, it is ambiguous whether negation has narrow or wide scope. 

\ea\label{NPI}
\ea\label{NPI1} 
  \glll      mɛ̀ɛ́ nyɛ́lɛ́ bùdì \\
          mɛ̀ɛ́ nyɛ̂-lɛ́ b-ùdì \\
              1{\SG}.{\PRS}.{\NEG} see-{\NEG} 2-person   \\
    \trans `I don't see any/the people.'
\ex\label{NPI2}    
  \glll      mɛ̀ɛ́ nyɛ́lɛ́ {\bfseries tɔ̀} bùdì \\
          mɛ̀ɛ́ nyɛ̂-lɛ́ tɔ̀ b-ùdì \\
              1{\SG}.{\PRS}.{\NEG} see-{\NEG} {\NPI} 2-person   \\
    \trans `I don't see any people.'
\z
\z

Gyeli can disambiguate narrow scope readings by using the negative polarity item ({\NPI}) {\itshape tɔ̀}, as in \REF{NPI2}, as is expected for the function of an {\NPI}: \citet[483]{Barker2018} shows that, ``an {\NPI} must always take narrow scope with respect to its licensing context.''

\subsection{Negative polarity}
\label{sec:gye:NegPol}

Gyeli has two negative polarity items (\textsc{npi}). One is {\itshape tɔ̀} and disambiguates scope readings under negation, as discussed in \sectref{sec:gye:Scope}. The other is {\itshape lìí} `yet', which is licensed under negation of the completive aspect, as shown in \REF{lii}. For further discussion, please see \citet[477]{Grimm2021}.
Both {\NPI}s are strictly limited to negative contexts. 

\ea \label{lii}
\ea \label{lii1}
  \glll  mɛ̀ pálɛ́ {\bfseries lìí} bâ \\
          mɛ pálɛ́ lìí bâ \\
          1\textsc{sg}.{\textsc{rec.}\PST} {\NEG}.{\PST} yet marry\\
    \trans `I am not yet married.'
\ex \label{lii2}
  \glll  mɛ̀ bâ mɔ̀ \\
          mɛ bâ mɔ̀ \\
          1\textsc{sg}.{\textsc{rec.}\PST} marry {\COMPL} \\
    \trans `I am already married.'    
\z
\z

In contrast to the opposition of English `still' and 
`anymore', Gyeli does not use an \textsc{npi} but the same lexical item {\itshape ná} `still' in affirmative and negated sentences, as shown in \REF{na} for both the standard and non-standard negation strategy in the present tense. In all cases, the continuous marker {\itshape ná} follows the inflected verb form, which is the auxiliary {\itshape tí} in \REF{na3}.  

\ea\label{na}
\ea\label{na1}
  \glll   líní á sílɛ́ dè mántúà á dyúwɔ́ {\bfseries ná} nzà \\
líní a-H sílɛ-H dè {H\backslash{}ma-ntúà} a-H dyúwɔ-H ná nzà  \\
when 1-\textsc{prs} finish-{\R} eat {{\OBJ}.{\LINK}\backslash{}6-mango} 1-{\PRS} feel-{\R} still $\emptyset$.9.hunger\\
    \trans `When s/he has eaten mangoes, s/he still feels hungry.' 
\ex\label{na2}
  \glll   líní á sílɛ́ dè mántúà àá dyúwɔ́lɛ́ {\bfseries ná} nzà \\
líní a-H sílɛ-H dè {H\backslash{}ma-ntúà} àá dyúwɔ-lɛ ná nzà  \\
when 1-\textsc{prs} finish-{\R} eat {{\OBJ}.{\LINK}\backslash{}6-mango} 1.{\PRS}.{\NEG} feel-{\NEG} still $\emptyset$.9.hunger\\
    \trans `When s/he has eaten mangoes, s/he does \textit{not} feel hungry anymore.'  %HC sanitycheck  
\ex\label{na3}
  \glll   líní á sílɛ́ dè mántúà à tí {\bfseries ná} dyúwɔ̀ nzà \\
líní a-H sílɛ-H dè {H\backslash{}ma-ntúà} a tí ná dyúwɔ nzà  \\
when 1-\textsc{prs} finish-{\R} eat {{\OBJ}.{\LINK}\backslash{}6-mango} 1 {\NEG} still feel $\emptyset$.9.hunger\\
    \trans `When s/he has eaten mangoes, s/he does not feel \textit{hungry} anymore.'
\z
\z

\subsection{Marking of {NP}s in the scope of negation}
There is no special marking of {\NP}s in the scope of negation except for the negative polarity item  {\itshape tɔ̀} which may appear in front of a NP to indicate narrow scope (\sectref{sec:gye:Scope}, \sectref{sec:gye:NegPol}). 

\subsection{Reinforcing negation}

Reinforcing negative replies (\sectref{sec:gye:NR}) is achieved prosodically through lengthening and increased volume. Also cleft constructions, as described in \sectref{sec:gye:Scope}, can reinforce negation.  Other means of negation reinforcement require future research. 

\subsection{Negation, coordination and complex clauses}
\label{sec:gye:Coord}

Coordination of affirmative and negated phrases seems generally to be avoided. Speakers use separate sentences or coordination with the opposition coordinator {\itshape ndí} `but' instead. Coordination of two negated phrases only requires negation marking on the first verb, while negation also has scope over the second coordinand, as in \REF{NCord}.

\ea\label{NCord} 
\glll mɛ̀ɛ́ gyímbɔ́lɛ́ nà dyà\\
mɛ̀ɛ́ gyímbɔ-lɛ nà dyà\\
1{\SG}.{\NEG}.{\PRS} dance-{\NEG} {\CONJ} sing\\
\trans `I don't dance and sing.'
\z

The negative polarity item {\itshape tɔ̀} (\sectref{sec:gye:NegPol}) is used as a negative coordinator. It can coordinate both {\NP}s, as in \REF{negcoord1}, and non-finite verbs in complex predicates, as in \REF{negcoord2}.

\ea\label{negcoord} 
\ea\label{negcoord1} 
  \glll mɛ̀ɛ́ kwálɛ̀lɛ̀ {\bfseries tɔ̀} bèfùmbí {\bfseries tɔ̀} mànjù\\
  mɛ̀ɛ́ kwàlɛ-lɛ tɔ̀ be-fùmbí tɔ̀ ma-njù\\
  1{\SG}.{\NEG}.{\PRS} like-{\NEG} {\NPI} 8-orange {\NPI} 6-banana\\ 
     \trans `I neither like oranges nor bananas.'
\ex\label{negcoord2} 
  \glll mɛ̀ɛ́ kwálɛ̀lɛ̀ {\bfseries tɔ̀} dyà {\bfseries tɔ̀} gyímbɔ̀\\
  mɛ̀ɛ́ kwàlɛ-lɛ tɔ̀ dyà tɔ̀ gyímbɔ\\
  1{\SG}.{\NEG}.{\PRS} like-{\NEG} {\NPI} sing {\NPI} dance\\ 
     \trans `I neither like to sing nor to dance.'
\z    
\z

The non-standard negation marker {\itshape tí} (\sectref{sec:gye:ti})  functions as a negative subordinator in infinitival subordinate clauses, as in \REF{tiSUBb}. This is different from simply negating the predicate of a dependent clause (\sectref{sec:gye:non-main}), as verb inflection for tense, aspect, mood, and polarity is generally lost in infinitival dependent clauses.  The affirmative counterpart in \REF{tiSUBa} either takes another subordinator such as the sequential marker {\itshape vɛ̀ɛ̀} or no special subordinator at all.

\ea\label{tiSUB}Context: Depiction of disease roaming in his body.
\ea\label{tiSUBa} 
  \glll  gbĩ́-gbĩ̀-gbĩ́-gbĩ̀-gbĩ́   à múà nà bábɛ̀ {\ob}({\bfseries vɛ̀ɛ̀}) wúmbɛ̀ wɛ̀{\cb} \\
            gbĩ́-gbĩ̀-gbĩ́-gbĩ̀-gbĩ́  a múà nà bábɛ̀ {\db}vɛ̀ɛ̀ wúmbɛ wɛ̀   \\
         {\IDEO}:roaming 1 {\PROSP} {\COM} $\emptyset$.7.illness {\db}{\textsc{seq}} want die \\
    \trans `He was about to be sick, wanting to die.'
\ex\label{tiSUBb} 
  \glll  gbĩ́-gbĩ̀-gbĩ́-gbĩ̀-gbĩ́   à múà nà bábɛ̀ {\ob}{\bfseries tí} wúmbɛ̀ wɛ̀{\cb} \\
            gbĩ́-gbĩ̀-gbĩ́-gbĩ̀-gbĩ́  a múà nà bábɛ̀ {\db}tí wúmbɛ wɛ̀   \\
         {\IDEO}:roaming 1 {\PROSP} {\COM} $\emptyset$.7.illness {\db}{\NEG} want die \\
    \trans `He was about to be sick, not wanting to die.'
\z
\z

The language does not have negative conjunctions such as {\itshape lest}. 

\section{Summary}\label{sec:gye:5}

Negation in Gyeli is mostly marked through syntactic means of negation auxiliaries, as shown in \tabref{tab:1:negmarkers}, which provides a summary of all negation markers in the language. 


\begin{table}\small
\caption{Gyeli negation markers}
\label{tab:1:negmarkers}
\begin{tabularx}{\textwidth}{XXl}
  \lsptoprule
Negation marker & Status & Distribution \\
\midrule
H...-\textit{lɛ} &  negation circumfix & present \\
\textit{sàlɛ́}/\textit{palɛ́} &  negation auxiliary & past tenses  \\
\textit{kálɛ̀} & negation auxiliary & future \\
\textit{dúù} `must not' & modal negation auxiliary & subjunctive, present \\
\textit{tí} &  negation auxiliary & imperative, infinitive, \\
& & present (predicate focus) \\
  \lspbottomrule
\end{tabularx}
\end{table}

In Gyeli clausal negation, both standard and non-standard negation strategies are used, as shown in \tabref{tab:1:negdist}.\footnote{Negation markers also take STAMP marker forms that are specific to the TAMP category they encode, as described in \sectref{sec:gye:SN}.} The distribution of standard versus non-standard negation correlates with a combination of specific tense-mood categories with certain clause types. There is, however, no entirely neat form-to-function mapping of negation markers to single tense-mood categories or clause types.  While past tenses and the future only occur with standard negation, present tense negation has several strategies, including standard and non-standard negation. Remarkably, the choice of different negation strategies coincides with different focus constructions, differentiating truth value focus and lexical verb focus.  Subjunctives and imperatives only have non-standard negation marking. Some clause types are associated with standard and non-standard negation, such as main, relative, and complement clauses. In contrast, non-verbal and existential clauses only take standard negation marking, while infinitival subordinate clauses and {\NP}s only allow non-standard negation. 

\begin{table}\small
\caption{Distribution of clausal negation markers}
\label{tab:1:negdist}
\begin{tabularx}{\textwidth}{Xll}
 \lsptoprule
Negation marker & {\TM} distribution & Clause types \\
\midrule
\multicolumn{3}{l}{\textbf{Standard negation}} \\
\midrule
{\itshape -lɛ} &  present & main, relative, complement,   \\
{\itshape sàlɛ́/palɛ́} &  past tenses &  non-verbal, existential\\
{\itshape kálɛ̀} &  future &  \\
\midrule
\multicolumn{3}{l}{\textbf{Non-standard clausal negation}} \\
\midrule
{\itshape dúù} `must not' & subjunctive, present & main, relative, complement \\
{\itshape tí} &  imperative, infinitive, & main, infinitival clause, {\NP}s\\
& present (predicate focus) & \\
\lspbottomrule
\end{tabularx}
\end{table}


While Gyeli has various strategies and markers for clausal negation, negative derivation is virtually absent. This ties in with the language's limited set of adjectives.  

\section*{Acknowledgements}
I would like to thank Ljuba Veselinova and Matti Miestamo for organizing the workshop on negation at 8\textsuperscript{th} {\itshape Conference on the Syntax of the World's Languages (SWL)} at Paris in 2018 and editing this book. I am grateful for the constructive comments given by two anonymous reviewers and feedback from the audience at SWL8, as well as to Héloïse Calame for very careful proofreading and typesetting.

% \clearpage
\section*{Abbreviations}

\noindent\begin{tabularx}{.45\textwidth}{lQ}
$\emptyset$ &  {\itshape $\emptyset$}-noun class\\
* &  ungrammatical form\\
1--9 &  agreement class 1--9\\
1\textsc{sg} &  first person singular \\
2\textsc{sg} &  second person singular\\
1\textsc{pl} &  first person plural\\
2\textsc{pl} &  second person plural\\
{\ADV} &  adverb\\
{\ANA} &  anaphoric marker\\
{\ATTR} &  attributive marker\\
%\itshape ba &  {\itshape ba}-noun class\\
%\itshape be &  {\itshape be}-noun class\\
{\COM} &  comitative marker\\
{\COMP} &  complementizer\\
{\COMPL} &  completive\\
{\CONJ} &  conjunction\\
{\COP} &  copula\\
{\DEM} &   demonstrative\\
{\FUT} &  future\\
H &  high tone\\
\scshape hab & habitual \\
HL &  falling contour tone\\
{\ID} &  identificational marker\\
{\IDEO} &  ideophone\\
\end{tabularx}
\begin{tabularx}{.45\textwidth}{lQ}
{\IMP} &    imperative\\
{\INCH} &  inchoative\\
L &   low tone\\
%\itshape le &  {\itshape le}-noun class\\
{\LOC} &   locative\\
%\itshape ma &  {\itshape ma}-noun class\\
{\N} &  noun\\
{\NCA} &  non-complete accomplishment\\
{\NEG} &  negation\\
{\NPI} &  negative polarity item\\
{\NP} &  noun phrase\\
{\OBJ}.{\LINK} &  object linking H tone\\
{\PL} & plural \\
{\PN} &  proper name\\
{\POSS} &  possessor pronoun\\
{\PREP} &  preposition\\
\scshape prf & perfect \\
{\PRIOR} &  priorative\\
{\PRO} &  pronoun\\
{\PROG} &  progressive\\
{\PROSP} &  prospective\\
{\PROX} &  proximal\\
\end{tabularx}

\begin{tabularx}{.45\textwidth}{lQ}
\textsc{prs} &  present\\ 
{\PST} &  past\\
{\Q} &  question marker\\
{\R} &  realis mood\\
\scshape rec &  recent \\
{\REL} &  relative clause\\
\scshape rem &  remote \\
{\RETRO} & retrospective \\
{\SBJV} &  subjunctive\\
\end{tabularx}
\begin{tabularx}{.45\textwidth}{lQ}
{\textsc{seq}} &  sequential marker\\
\scshape sg &  singular\\
\textsc{stamp} &  subject agreement, tense, aspect, mood, polarity clitic\\
{\SUB} &  subordinate\\
{\TM} &  tense-mood\\
{\V} &  verb\\
\\
\end{tabularx}

\printbibliography[heading=subbibliography,notkeyword=this]
\end{document}
